\SourceLeaf{422}\noindent Anstatt aus den vier Gleichungen, welche die Substitution dieser Werthe in den beiden Gleichungen (1) und (2) Art. 2 ergibt, die unbekannten Elemente $P,Q^0,R^0,S^0$ durch Elimination abzuleiten, gelangen wir zu demselben Ziele viel leichter durch Benutzung des Calculs der imaginären Grössen, indem wir in Folge der Formel (4) Art. 2 von den beiden Gleichungen
\[
\begin{aligned}
1&=\frac{s'S^0}{t't'}e^{i(P+p'-A')}-r'R^0e^{i(Q^0+q'-B')}\\
1&=\frac{s''S^0}{t''t''}e^{i(P+p''-A'')}-r''R^0e^{i(Q^0+q''-B'')}
\end{aligned}
\]
ausgehen, die sich, wenn wir zur Abkürzung
\[
\begin{aligned}
\frac{s'}{t't'}e^{i(p'-A')}&=a'\\
\frac{s''}{t''t''}e^{i(p''-A'')}&=a''\\
r'e^{i(q'-B')}&=b'\\
r''e^{i(q''-B'')}&=b''\\
S^0e^{iP}&=x\\
R^0e^{iQ^0}&=y
\end{aligned}
\]
setzen, in folgende verwandeln
\[
\begin{aligned}
1&=a'x-b'y\\
1&=a''x-b''y
\end{aligned}
\]
woraus man
\[
\begin{aligned}
x&=\frac{b''-b'}{a'b''-a''b'}=\frac{\dfrac{b''}{b'}-1}{\dfrac{b''}{b'}-\dfrac{a''}{a'}}.\frac{1}{a'}\\[1ex]
y&=\frac{a''-a'}{a'b''-a''b'}=\frac{\dfrac{a''}{a'}-1}{\dfrac{b''}{b'}-\dfrac{a''}{a'}}.\frac{1}{b'}
\end{aligned}
\]
erhält. Es ergeben sich hieraus folgende entwickelte Rechnungsvorschriften. Man setze (I)
\SourceLeaf{423}
\[
\begin{aligned}
\frac{t't'}{s'}.\frac{s''}{t''t''}.\cos(A'-A''-p'+p'')&=\mathfrak{A}\\
\frac{t't'}{s'}.\frac{s''}{t''t''}.\sin(A'-A''-p'+p'')&=\mathfrak{A}_1\\
\frac{r''}{r'}.\cos(B'-B''-q'+q'')&=\mathfrak{B}\\
\frac{r''}{r'}.\sin(B'-B''-q'+q'')&=\mathfrak{B}_1
\end{aligned}
\]
wodurch also
\[
\begin{aligned}
\frac{a''}{a'}&=\mathfrak{A}+i\mathfrak{A}_1\\
\frac{b''}{b'}&=\mathfrak{B}+i\mathfrak{B}_1
\end{aligned}
\]
wird. Man bestimme ferner die sechs Grössen $u,U,v,V,w,W$ aus den Gleichungen (II)
\[
\begin{aligned}
\mathfrak{A}-1&=u\cos U\\
\mathfrak{A}_1&=u\sin U\\
\mathfrak{B}-1&=v\cos V\\
\mathfrak{B}_1&=v\sin V\\
\mathfrak{B}-\mathfrak{A}&=w\cos W\\
\mathfrak{B}_1-\mathfrak{A}_1&=w\sin W
\end{aligned}
\]
und zwar so, dass $u,v,w$ positiv werden. Es wird dann
\[
\begin{aligned}
\frac{b''}{b'}-1&=ve^{iV}\\
\frac{a''}{a'}-1&=ue^{iU}\\
\frac{b''}{b'}-\frac{a''}{a'}&=we^{iW}
\end{aligned}
\]
und folglich
\[
\begin{aligned}
x&=\frac{t't'}{s'}.\frac{v}{w}e^{i(V-W+A'-p')}\\
y&=\frac{u}{r'w}e^{i(U-W+B'-q')}
\end{aligned}
\]
woraus man leicht schliesst, dass (III)
\SourceLeaf{424}
\[
\begin{aligned}
P&=V-W+A'-p'\\
Q^0&=U-W+B'-q'\\
R^0&=\frac{u}{r'w}\\
S^0&=\frac{t't'}{s'}.\frac{v}{w}
\end{aligned}
\]
Die vierzehn Formeln I, II, III enthalten die vollständige möglich einfachste Auflösung unsrer Aufgabe.

Es verdient noch bemerkt zu werden, dass für $r'=r''$, (sei es, dass die vergleichenden Beobachtungen diese Gleichheit ergeben, oder dass man in Ermangelung solcher Beobachtungen die Veränderlichkeit von $R$ während der beiden Beobachtungssätze zu berücksichtigen nicht im Stande ist)
\[
\begin{aligned}
V&=\tfrac12(B'-B''-q'+q'')\pm90^0\\
v&=\pm2\sin\tfrac12(B'-B''-q'+q'')
\end{aligned}
\]
wird, wo die obern oder die untern Zeichen gelten, je nachdem
\[
\sin\tfrac12(B'-B''-q'+q'')
\]
positiv oder negativ ist.

\Article{8}
Zur Erläuterung dieser Vorschriften fügen wir noch die vollständige Berechnung eines Beispiels bei. Die Rechnung ist mit siebenzifrigen Logarithmen geführt, also viel schärfer, als für die Ausübung nöthig ist, wo fünfzifrige Logarithmen immer zureichen.

Am 24. März 1841 wurde die Schwingungsdauer des Bifilarmagnetometers aus Beobachtungen, welche $1^{St}21'$ umfassten (wie sich von selbst versteht, nach gehöriger Reduction auf unendlich kleine Schwingungen) $=28''89071=t'$ gefunden; die Stellung der Spiegelalhidade war $154^020'30''=A'$, die der Alhidade des Schiffchens $=27^040'25''=B'$. Im Mittel aus mehreren über jenen Zeitraum gleichförmig vertheilten Bestimmungen war der Stand $994.33$ Skalentheile, also da der Lothfaden der Skalenzahl $1000$ entspricht, und ein Skalentheil $21''5835$ beträgt, $p'=-2'2''38$. Aus ganz gleichzeitigen Beobachtungen fand sich die Schwingungsdauer des Unifilarmagnetometers im magnetischen
\SourceLeaf{425}Observatorium $=20''72725$, und der Stand im Mittel $=881.80$ Skalentheile. Als Normalstand wurde der mittlere Stand aus den täglichen Aufzeichnungen im Februar $888.40$ gewählt (welchem übrigens die absolute Declination $18^011'54''$ entspricht); da ein Skalentheil am Unifilarmagnetometer $21''3489$ beträgt, so findet sich daraus $q'=-2'20''90$.

Der mittlere Thermometerstand (aus Aufzeichnungen unmittelbar vor dem Anfange und gleich nach dem Schluss der Beobachtungen) war im Kasten des Bifilarmagnetometers $+6^096$, bei der obern Befestigung der Aufhängungsdrähte $+7^06$, im Kasten des Unifilarmagnetometers $+7^045$, alles nach Réaumur.

Auf gleiche Weise war für einen zweiten Satz von Beobachtungen am folgenden Tage
\[
\begin{aligned}
t''&=108''17\\
A''&=151^027'30''\\
B''&=185\;59\;35\\
p''&=-24'33''07\\
q''&=+2\;42.04
\end{aligned}
\]
die Schwingungsdauer des Unifilarmagnetometers $=20''73117$, die Thermometerstände in derselben Ordnung wie oben $+6^036$, $+7^00$, $+7^01$.

Zur Berechnung des Einflusses der Temperatur setze ich
\[
\alpha=0.000016,\quad\beta=0.000024,\quad\gamma=\gamma'=0.000765
\]
den letztern Werth nach \textsc{Hansteen}, da eigne entscheidende Bestimmungen zur Zeit noch fehlen. Es folgt hieraus nach den Formeln (1) und (3) des 5. Art., wenn wir $20''72=\theta^0$ zur Normalschwingungsdauer des Unifilarmagnetometers wählen,
\[
\begin{aligned}
\log r'&=-0.0001376\\
\log s'&=-0.0000043\\
\log r''&=-0.0002155\\
\log s''&=-0.0000044
\end{aligned}
\]
In Folge der abgekürzten Formeln (2) und (4) a. a. O. würde man setzen können
\[
\begin{aligned}
\log s&=-\alpha k(c''-c)\\
\log r&=2\log\frac{\theta^0}{\theta}+(\gamma+2\alpha)k(c'-c)-\alpha k(c''-c)
\end{aligned}
\]
\SourceLeaf{426}wenn $k$ den Modulus der briggischen Logarithmen bezeichnet, also mit obigen Werthen von $\alpha,\beta,\gamma$
\[
\begin{aligned}
\log s&=0.00000695(c''-c)\\
\log r&=2\log\frac{\theta^0}{\theta}+0.0003461(c'-c)-0.00000695(c''-c)
\end{aligned}
\]
woraus für unsre Beobachtungen folgt
\[
\begin{aligned}
\log r'&=-0.0001386\\
\log s'&=-0.0000044\\
\log r''&=-0.0002169\\
\log s''&=-0.0000044
\end{aligned}
\]
also kaum merklich von obigen Werthen verschieden.

Nach diesen Vorbereitungen sind die Hauptmomente der Rechnung selbst folgende:
\[
\begin{aligned}
A'-A''-p'+p''&=+2^030'29''31\\
B'-B''-q'+q''&=-158\;14\;7.06\\
\log\frac{t't's''}{s't''t''}&=8.8524525\\
\log\frac{r''}{r'}&=9.9999221
\end{aligned}
\]
Hieraus nach I
\[
\begin{aligned}
\log\mathfrak{A}&=8.8520363\\
\log\mathfrak{A}_1&=7.4935432\\
\log\mathfrak{B}&=9.9678043\,\mathrm{n}\\
\log\mathfrak{B}_1&=9.5690569\,\mathrm{n}
\end{aligned}
\]
woraus ferner
\[
\begin{aligned}
\log(\mathfrak{A}-1)&=9.9679562\,\mathrm{n}\\
\log(\mathfrak{B}-1)&=0.2852304\,\mathrm{n}\\
\log(\mathfrak{B}-\mathfrak{A})&=9.9998589\,\mathrm{n}\\
\log(\mathfrak{B}_1-\mathfrak{A}_1)&=9.5726915\,\mathrm{n}
\end{aligned}
\]
Hienach ergeben die Formeln II
\SourceLeaf{427}
\[
\begin{aligned}
U&=179^048'28''15\\
V&=190\;52\;52.91\\
W&=200\;30\;14.79\\
\log u&=9.9679586\\
\log v&=0.2931101\\
\log w&=0.0282829
\end{aligned}
\]
und endlich die Formeln III
\[
\begin{aligned}
P^0&=144^045'10''50\\
Q^0&=7\;0\;59.26\\
\log R^0&=9.9398133\\
\log S^0&=3.1863476
\end{aligned}
\]

\Article{9}
Noch mehr lässt sich die Aufgabe generalisiren, indem man \textit{vier verschiedene} Beobachtungssätze zum Grunde legt, zwei für den Stand, zwei für die Schwingungsdauer, wobei man zugleich die Voraussetzung fahren lässt, dass diese und jene beziehungsweise denselben Werthen von $B$ entsprechen. Man hat dabei zwar den Vortheil, die Beobachtungen für den Stand des Bifilarmagnetometers nach dem in den \textit{Resultaten} für 1836. II. [S. unten] beschriebenen Verfahren bei einem beinahe ganz beruhigten Zustande des Magnetstabes machen zu können: allein dieser Vortheil verliert seinen Werth durch den Umstand, dass man genöthigt bleibt, für alle vier Sätze am Hülfsmagnetometer Schwingungsdauer und Stand zugleich zu beobachten, also letztern doch aus Elongationen bestimmen muss. Es erhellt also, dass diese Methode doppelt so viele Arbeit verursacht, als die des 7. Art., welche ausserdem den Vorzug einer so sehr einfachen Berechnung hat, während die directe Bestimmung der Elemente aus vier getrennten Beobachtungssätzen bei weitem weitläufiger ausfällt, daher wir auch ihre in mathematischer Beziehung nicht uninteressante Entwickelung lieber auf einen andern Ort versparen.

\Article{10}
Es verdient noch bemerkt zu werden, dass wenn man bei der Bestimmung des Standes aus beobachteten Elongationen das [S. unten] \textit{Resultate} 1836 II
\SourceLeaf{428}angezeigte Verfahren schlechthin anwendet, die ungleiche Geltung der Skalentheile in Bogentheilen einen Fehler erzeugt; der desto grösser ist, je weiter der Stand von der Mitte der Skale abliegt. Verlangt man also ganz scharfe Resultate, so muss man jenes Verfahren nicht unmittelbar auf die in den Elongationen abgelesenen Skalentheile, sondern auf die nach strenger Formel in Bogentheile verwandelten Abstände der Elongationen von der Mitte der Skale anwenden. Ist der Stand nahe bei der Mitte, so ist allerdings jener Fehler unerheblich, und man wird daher immer die Stellung des Spiegels oder den Werth von $A$ so wählen, dass der Stand von der Mitte wenig abweiche, oder dass $p$ klein werde. Bei der ersten Bestimmung der Elemente ist dies freilich nur durch einen vorläufigen Versuch (auf ähnliche Art wie im Art. 6) zu erreichen: besitzt man aber schon eine genäherte Kenntniss der Elemente $P,Q,R$, so wird man zu diesem Zweck lieber eine Rechnung anwenden, welcher man am bequemsten folgende (aus Art. 2. Formel (1) oder (4) leicht abzuleitende) Gestalt gibt. Man bestimme einen Winkel $\psi$ durch die Formel
\[
\operatorname{tang}\psi=\frac{1-R}{1+R}.\operatorname{tang}\tfrac12(Q-B)=\operatorname{tang}\tfrac12\varphi^2.\operatorname{tang}\tfrac12(Q-B)
\]
und zwar so, dass $\psi$ in demselben Quadranten gewählt wird, in welchem $\tfrac12(Q-B)$ liegt, und setze dann
\[
A=\psi+P-\tfrac12(Q-B)
\]

\Article{11}
Es bleibt nun noch übrig, den Zusammenhang zu entwickeln, in welchem die Beobachtungen am Bifilarmagnetometer in der Transversalstellung mit den Veränderungen der Elemente stehen.

Es wird vorausgesetzt, dass die nach der Vorschrift von Art. 3 bestimmte Transversalstellung sich auf die Normalwerthe der Elemente beziehe: das Schiffchen ist also so gestellt, dass beim Ruhestande die magnetische Axe des Magnetstabes einen rechten Winkel mit dem magnetischen Normalmeridian macht, wenn das Verhältniss der magnetischen Richtungskraft zur statischen wie $R^0$ zu $1$ ist; der Spiegel hingegen so, dass bei jener Stellung das Bild des durch den Lothfaden bezeichneten Skalenpunkts auf dem Fadenkreuze des Beobachtungsfernrohrs erscheint. Es ist also, wenn wir die unter jenen Umständen Statt findende Schwingungsdauer mit $t^0$ bezeichnen,
\SourceLeaf{429}
\[
\begin{aligned}
\sin\varphi&=R^0\\
A&=P\mp\varphi\\
B&=Q^0\mp(90^0+\varphi)\\
\frac{S^0}{t^0t^0}&=\cos\varphi
\end{aligned}
\]
wo die doppelten Zeichen sich auf die westliche oder östliche Stellung des Nordpols des Magnetstabes beziehen. Indem wir nun die Zeichen
\[
p,\quad Q=Q^0+q,\quad R=rR^0,\quad S=sS^0,\quad t
\]
in der bisherigen allgemeinen Bedeutung beibehalten, geben die Formeln (1) und (6) des 2. Art.
\[
\begin{aligned}
\sin(\varphi\pm p)&=rR^0\cos(p-q)\\
\frac{sS^0}{tt}\sin(\varphi\pm p)&=rR^0\cos(\varphi\pm q)
\end{aligned}
\]
oder
\[
\begin{aligned}
r&=\frac{\sin(\varphi\pm p)}{\sin\varphi\cos(p-q)}\quad .\ .\ .\ .\ .\ .\ .\ .\ .\quad(1)\\
\frac{st^0t^0}{tt}&=\frac{\cos(\varphi\pm q)}{\cos\varphi\cos(p-q)}\quad .\ .\ .\ .\ .\ .\ .\ .\ .\quad(2)
\end{aligned}
\]

\Article{12}
Die wichtigste Anwendung des Bifilarmagnetometers ist die Bestimmung der Veränderungen der horizontalen Intensität, mit welchen die Veränderungen von $R$ durch die oben (Art. 5) gegebene Formel
\[
r=\frac{f^0g^0hMT}{fgh^0M^0T^0}
\]
zusammenhängen. Man muss sich hiebei erinnern, dass $T^0$ die Anfangs gewählte Normalintensität, $M^0$ das auf die Normaltemperatur reducirte magnetische Moment des Magnetstabes nach dessen magnetischem Zustande zur Zeit der Bestimmung der Constanten ausdrückt. Bezeichnen wir das eben so auf die Normaltemperatur reducirte magnetische Moment für eine unbestimmte Zeit mit $\mathfrak{M}$, und setzen
\[
\frac{\mathfrak{M}}{M^0}=\lambda,\qquad\frac{T^0}{\lambda}=\mathfrak{T}
\]
so wird unter den im 4. Art. besprochenen Bedingungen $\lambda$ ein für eine mässige
\SourceLeaf{430}Zeit, z. B. für Einen Tag, wie constant zu betrachtender Coëfficient sein, und so wie dieser zugleich mit $\mathfrak{M}$ allmählig sehr langsam abnimmt, wird $\mathfrak{T}$ allmählig zunehmen und stets diejenige horizontale Intensität ausdrücken, bei welcher unter der Normaltemperatur das Verhältniss der magnetischen und der statischen Richtungskraft dem Verhältnisse $R^0:1$ gleich wird. Da nun obige Gleichung die Form
\[
T=\frac{\mathfrak{M}fgh^0}{Mf^0g^0h}.r\mathfrak{T}
\]
annimmt, wo der erste Factor $\frac{\mathfrak{M}fgh^0}{Mf^0g^0h}$ bloss von der Temperatur abhängt, und (wenn wir die Bezeichnungen des 5. Art. beibehalten) durch
\[
1+(\gamma+2\beta-\alpha)c+(\theta-\tfrac12\alpha)(c''-c)
\]
ausgedrückt werden kann; $r$ hingegen durch combinirte gleichzeitige Standbeobachtungen am Bifilarmagnetometer in der transversalen Stellung für $p$, und am Unifilarmagnetometer für $q$, nach Formel (1) des vorhergehenden Art. für jeden Augenblick bestimmbar ist: so erhellt, dass sich auf diese Weise die Veränderungen der Intensität in den kleinsten Zeitfristen mit grösster Schärfe verfolgen lassen, so lange es nur darauf ankommt, die veränderten Intensitäten während eines mässigen Zeitraumes, z. B. während eines vierundzwanzigstündigen Termins, oder während der zu einer absoluten Intensitätsbestimmung vermittelst des Unifilarmagnetometers erforderlichen Zeit, \textit{unter sich} zu vergleichen. Indem man bei einer solchen absoluten Intensitätsbestimmung zu den Reductionen der einzelnen Operationen auf einerlei Normalintensität (vergl. \textit{Intensitas vis magneticae} Art. 10 und 22) die gleichzeitigen Beobachtungen am Bifilarmagnetometer verwendet (was auch an sich vortheilhafter ist, als der a. a. O. empfohlene Gebrauch eines zweiten Unifilarmagnetometers), erhält man zugleich die Kenntniss des für diese Zeit gültigen Werths von $\mathfrak{T}$ in absolutem Maasse. Wenn man nun solche absolute Bestimmungen von Zeit zu Zeit wiederholt, so bleibt man fortwährend von den etwaigen allmähligen Veränderungen von $\mathfrak{T}$ in Kenntniss, und kann dieselben für die Zwischenzeit nach geometrischer Progression durch Interpolation ohne merklichen Fehler ansetzen, und sonach sämmtliche Veränderungen der Intensität nach allen ihren Abwechslungen in absolutem Maasse angeben. Übrigens versteht sich von selbst, dass, wenn nach längerer Zwischenzeit, in Folge der Säcularänderungen der magnetischen Declination und horizontalen In%
\SourceLeaf{431}tensität, oder beträchtlicher Schwächung des Stabmagnetismus, $p$ und $q$ aufhören innerhalb mässiger Grenzen zu bleiben (wozu aber die Fälle grosser ausserordentlicher Anomalien nicht gerechnet werden müssen), man eine zweckmässige Abänderung an der Stellung des Schiffchens, des Spiegels, und wenn man es rathsam findet auch des Abstandes der Aufhängungsdrähte vornehmen, und so eine neue Reihe von Beobachtungen mit veränderten Elementen anfangen wird.

\Article{13}
So lange $p$ und $q$ nur klein sind, wird man für alle Zwecke, wo die grösste Schärfe nicht gefordert wird, anstatt der strengen Formel (1) eine abgekürzte anwenden können, wo $q$ ganz herausfällt, also gleichzeitige Beobachtungen am Unifilarmagnetometer gar nicht gebraucht werden: dies gilt namentlich von den gewöhnlichen Terminsbeobachtungen. Anstatt jener Formel kann man nemlich setzen
\[
r=1\pm\operatorname{cotang}\varphi.\operatorname{tang}p\quad\text{oder auch}\quad r=1\pm\tfrac12\operatorname{cotang}\varphi.\operatorname{tang}2p
\]
Da nun, wenn $n$ den Unterschied des abgelesenen Skalentheils von demjenigen, auf welchen der Lothfaden sich bezieht, und $d$ die horizontale Entfernung der Mitte des Spiegels von letzterm Punkte in Skalentheilen gemessen, bedeutet,
\[
\operatorname{tang}2p=\frac{n}{d}
\]
ist, so verwandelt sich diese Formel in
\[
r=1\pm\frac{n}{2\operatorname{tang}\varphi.d}
\]
und es wird dann zugleich, hinlänglich genau,
\[
T=\mathfrak{T}(1\pm\frac{n}{2\operatorname{tang}\varphi.d}+(\gamma+2\beta-\alpha)c)
\]
wenn man das geringfügige Glied $(\beta-\tfrac12\alpha)(c''-c)$ weglässt. Bei dem hiesigen Apparate ist $d=4778{,}3$ Millimeter, und nach den Resultaten der im 8. Art. als Beispiel geführten Rechnung ergibt sich $\varphi=60^031'37''9$, also $2d\operatorname{tang}\varphi=16910$. Mit den daselbst gebrauchten Werthen von $\alpha,\beta,\gamma$ erhält man also
\[
T=(1+\frac{n+13{,}65c}{16910})\mathfrak{T}
\]
